% =============================================================================
% Experimental Protocol: gen_2 Neural Track Extrapolation for LHCb Allen HLT1
%
% Author: G. Scriven
% Date: April 2026
% =============================================================================
\documentclass[a4paper,11pt]{article}
\usepackage[margin=2.5cm]{geometry}

% --- Packages ---
\usepackage{amsmath,amssymb}
\usepackage{graphicx}
\usepackage{booktabs}
\usepackage{hyperref}
\usepackage{xcolor}
\usepackage{algorithm}
%\usepackage{algpseudocode}  % not available
%\usepackage{multirow}  % not available
\usepackage{array}
\usepackage{bm}
\usepackage{microtype}          % improves line breaking / reduces overfull boxes

% Custom commands — SI unit replacements (siunitx not available)
\newcommand{\SI}[2]{#1~#2}
\newcommand{\SIrange}[3]{#1--#2~#3}
\newcommand{\num}[1]{#1}
\newcommand{\si}[1]{#1}
\newcommand{\kilo}{k}
\newcommand{\byte}{B}
\newcommand{\tesla}{T}
\newcommand{\mega}{M}
\newcommand{\hertz}{Hz}
\newcommand{\giga}{G}
\newcommand{\eV}{eV}
\newcommand{\mm}{mm}
\newcommand{\minute}{min}
\newcommand{\qop}{q/p}
\newcommand{\dz}{\Delta z}
\newcommand{\Bfield}{\mathbf{B}}
\newcommand{\state}{\mathbf{y}}
\newcommand{\Jac}{\mathbf{J}}
\newcommand{\loss}{\mathcal{L}}
\newcommand{\btheta}{\bm{\theta}}

\begin{document}

% ============================================================================
\title{Experimental Protocol: Second-Generation Neural Track Extrapolation\\
for the LHCb Allen HLT1 Trigger}

\author{G.~Scriven\\\small Maastricht University, UHasselt, Nikhef}

\date{\today}

\begin{abstract}
This document specifies the experimental protocol for the \texttt{gen\_2}
study of neural network--based track extrapolation in the LHCb experiment.
The study trains multi-layer perceptrons (MLPs), physics-informed neural
networks (PINNs), and Runge-Kutta--structured PINNs (RK\_PINNs) to
approximate the fourth-order Runge-Kutta (RK4) numerical extrapolator that
propagates charged-particle states through the LHCb dipole magnetic field.
Nine model configurations are trained on \num{50e6} RK4-generated tracks
spanning $\dz \in [25, 10000]$~mm and $p \in [1, 200]$~GeV/$c$.  The
protocol addresses two research questions: (RQ1) whether GPU-deployable
models ($\leq$~\SI{64}{\kilo\byte}) achieve accuracy competitive with RK4
for integration into Allen HLT1, and (RQ2) whether physics-informed
training improves accuracy-per-parameter relative to purely data-driven
training.  All data generation parameters, model architectures, loss
functions, training hyperparameters, evaluation metrics, known confounds,
and reproducibility provisions are specified in full.
\end{abstract}

\maketitle

% ============================================================================
\section{Introduction}
\label{sec:introduction}

The LHCb detector~\cite{LHCb_detector} reconstructs the trajectories of
charged particles produced in proton--proton collisions at the Large Hadron
Collider.  Track reconstruction requires propagating particle states
through the inhomogeneous LHCb dipole magnetic field from one detector
element to another---a procedure called \emph{track extrapolation}.  The
production extrapolator is a fourth-order Runge-Kutta (RK4)
integrator~\cite{RK4_extrapolator} that evaluates the Lorentz force
equation at each integration step, with the magnetic field obtained by
trilinear interpolation of a measured three-dimensional field
map~\cite{LHCb_field_map}.

In the upgraded LHCb detector operating at Run~3 and beyond, the
first-level software trigger (HLT1) runs entirely on
GPUs~\cite{Allen_GPU}.  The Allen framework executes real-time pattern
recognition at the full \SI{30}{\mega\hertz} bunch-crossing rate.  Track
extrapolation---called $\mathcal{O}(10^6)$ times per event---is a
significant contributor to the per-event processing time.  Replacing the
RK4 integrator with a neural network that can be evaluated efficiently in
GPU registers or constant memory offers a potential throughput improvement,
provided the accuracy degradation is acceptable for the downstream Kalman
filter~\cite{Kalman_filter}.

This document defines the protocol for \texttt{gen\_2}, the second
generation of this investigation.  The first generation
(\texttt{gen\_1})~established feasibility: a two-layer MLP with
architecture $[64, 64]$ achieved a mean position error of \SI{0.49}{\mm}
and was deployable within the \SI{64}{\kilo\byte} CUDA \texttt{\_\_constant\_\_}
memory limit.  However, \texttt{gen\_1} had two shortcomings: (i)~training
data used $\dz_\text{min} = \SI{100}{\mm}$, which does not cover the
\SIrange{25}{30}{\mm} spacing between VELO silicon pixel sensors, and
(ii)~the larger models that achieved better accuracy exceeded the
\SI{64}{\kilo\byte} deployment constraint.

\texttt{gen\_2} addresses both issues by extending the $\dz$ range down
to \SI{25}{\mm} and constraining all model architectures to
$\leq$~\SI{64}{\kilo\byte}.  It further introduces physics-informed
training as a mechanism to improve accuracy without increasing model size.

\subsection{Research Questions}

\noindent\textbf{RQ1 --- Engineering viability.}
Are GPU-deployable models ($\leq$~\SI{64}{\kilo\byte}) accurate enough for
integration into the Allen HLT1 workflow?  Specifically, does the MLP
extrapolator achieve position and slope accuracy competitive with the RK4
ground truth, particularly at the short $\dz$ values
(\SIrange{25}{100}{\mm}) relevant to VELO track fitting?

\noindent\textbf{RQ2 --- Physics-informed training.}
For a given architecture size, does embedding the Lorentz force ordinary
differential equation (ODE) as a soft constraint during training (PINN,
RK\_PINN) yield a model that generalises better---measured by
accuracy-per-parameter---compared to purely data-driven (MLP) training?

% ============================================================================
\section{Data Generation Protocol}
\label{sec:data_generation}

\subsection{Ground-Truth Integrator}

Training data are generated by a fixed-step-size RK4 integrator
(\texttt{experiments/\allowbreak gen\_2/\allowbreak utils/\allowbreak rk4\_propagator.py}) that propagates the
five-dimensional state vector
$\state = (x, y, t_x, t_y, q/p)$
through the LHCb magnetic field from an initial longitudinal position
$z_0$ to $z_f = z_0 + \dz$.  The integration step size is
$h = \SI{5}{\mm}$, chosen to be well within the stability region of RK4
for the field gradients encountered in the LHCb tracking
volume~\cite{RK4_extrapolator}.  Momentum ($q/p$) is conserved along the
trajectory (no material interactions).

\subsection{Magnetic Field Map}

The magnetic field is evaluated by trilinear interpolation of the measured
field map \texttt{twodip.rtf}, which tabulates $B_x$, $B_y$, $B_z$ on an
$81 \times 81 \times 146$ Cartesian grid:
\begin{itemize}
  \item $x, y \in [-4000, 4000]$~mm at \SI{100}{\mm} spacing,
  \item $z \in [-500, 14000]$~mm with variable spacing,
  \item peak $|B_y| \approx \SI{1.03}{\tesla}$ at $z \approx \SI{5007}{\mm}$.
\end{itemize}
The same interpolation module
(\texttt{experiments/gen\_2/utils/magnetic\_field.py}) is used for data
generation (NumPy backend) and for physics-informed training (PyTorch
backend with \texttt{torch.nn.functional.grid\_sample}), ensuring
consistency between the ground truth and the PDE residual computation.

\subsection{Sampling Strategy}

Each training sample consists of a randomly drawn initial state and step
size.  The parameter ranges and sampling distributions are:
\begin{align}
  z_0        &\sim \mathcal{U}(0,\; 14000 - \dz) \text{ mm}, \label{eq:z0} \\
  \dz        &\sim \mathcal{U}(25,\; 10000) \text{ mm}, \label{eq:dz} \\
  p          &\sim \text{LogU}(1,\; 200) \text{ GeV}/c, \label{eq:p} \\
  x_0        &\sim \mathcal{U}(-300,\; 300) \text{ mm}, \label{eq:x0} \\
  y_0        &\sim \mathcal{U}(-250,\; 250) \text{ mm}, \label{eq:y0} \\
  t_{x,0}    &\sim \mathcal{U}(-0.3,\; 0.3), \label{eq:tx0} \\
  t_{y,0}    &\sim \mathcal{U}(-0.25,\; 0.25), \label{eq:ty0}
\end{align}
where $\mathcal{U}$ denotes uniform sampling and $\text{LogU}$ denotes
log-uniform (uniform in $\log p$).  The charge sign is fixed to $q = -1$
(MagDown polarity).  Tracks that diverge beyond $|x|$ or $|y| >
\SI{5000}{\mm}$, or whose slopes exceed $|t_x|$ or $|t_y| > 2.0$, are
rejected as unphysical.

The extension of $\dz_\text{min}$ from \SI{100}{\mm} (\texttt{gen\_1}) to
\SI{25}{\mm} is motivated by the VELO sensor pitch: the Kalman filter
propagates states between consecutive VELO planes separated by
\SIrange{25}{30}{\mm}, and the neural extrapolator must handle these
short steps.

\subsection{Input and Output Representation}
\label{sec:io_format}

The neural network input is a six-dimensional vector:
\begin{equation}
  \mathbf{x} = (x_0,\; y_0,\; t_{x,0},\; t_{y,0},\; q/p,\; \dz),
  \label{eq:input}
\end{equation}
with units $(\text{mm},\; \text{mm},\; -,\; -,\; \text{MeV}^{-1},\;
\text{mm})$.  The output is the four-dimensional extrapolated state:
\begin{equation}
  \mathbf{y} = (x_f,\; y_f,\; t_{x,f},\; t_{y,f}),
  \label{eq:output}
\end{equation}
omitting $q/p$ (conserved) and $z_f$ (deterministic: $z_f = z_0 + \dz$).

\subsection{Dataset Construction}

A total of \num{50e6} tracks are generated in 2000 parallel HTCondor
batches of \num{25000} tracks each
(\texttt{condor/submit\_datagen.sub}).  Each batch uses a
deterministic seed:
\begin{equation}
  s_i = 42 + i \times 7919, \quad i = 0, 1, \ldots, 1999,
  \label{eq:seed}
\end{equation}
where 7919 is a large prime chosen to space seeds far apart in the
pseudorandom number generator's period.  After all batches complete, the
individual \texttt{.npz} files are merged
(\texttt{data/merge\_batches.py}) into a single dataset
\texttt{train\_50M\_dz25.npz} containing arrays:
\begin{itemize}
  \item $\mathbf{X} \in \mathbb{R}^{N \times 6}$: input states,
  \item $\mathbf{Y} \in \mathbb{R}^{N \times 4}$: output states,
  \item $\mathbf{P} \in \mathbb{R}^{N}$: momenta (GeV/$c$),
  \item $\mathbf{Z} \in \mathbb{R}^{N}$: initial $z_0$ positions.
\end{itemize}
The dataset is split by random permutation (seed~42) into:
\begin{center}
\begin{tabular}{lrr}
\toprule
Split & Fraction & Samples \\
\midrule
Train      & 80\% & \num{40e6} \\
Validation & 10\% & \num{5e6}  \\
Test       & 10\% & \num{5e6}  \\
\bottomrule
\end{tabular}
\end{center}

% ============================================================================
\section{Physics: Equations of Motion}
\label{sec:physics}

A charged particle traversing the LHCb magnetic field $\Bfield =
(B_x, B_y, B_z)$ follows the Lorentz force law.  In the LHCb
$z$-parameterisation~\cite{RK4_extrapolator}, where $z$ (the beam axis)
serves as the independent variable, the equations of motion are:
\begin{align}
  \frac{dx}{dz}   &= t_x, \label{eq:dx} \\
  \frac{dy}{dz}   &= t_y, \label{eq:dy} \\
  \frac{dt_x}{dz} &= \kappa \, \mathcal{N}
    \bigl[ t_x t_y B_x - (1 + t_x^2) B_y + t_y B_z \bigr],
    \label{eq:dtx} \\
  \frac{dt_y}{dz} &= \kappa \, \mathcal{N}
    \bigl[ (1 + t_y^2) B_x - t_x t_y B_y - t_x B_z \bigr],
    \label{eq:dty}
\end{align}
where
\begin{equation}
  \kappa = \frac{q}{p} \times c_\text{light},
  \qquad
  \mathcal{N} = \sqrt{1 + t_x^2 + t_y^2},
  \label{eq:kappa}
\end{equation}
and $c_\text{light} = 2.99792458 \times 10^{-4}$~mm/(MeV\,$\cdot$\,T)
is the conversion factor from SI to LHCb internal units.  The dominant
field component is $B_y$ (vertical), which causes horizontal ($x$)
bending.  The slopes $t_x = dx/dz$ and $t_y = dy/dz$ are dimensionless
direction tangents.

Equations~\eqref{eq:dx}--\eqref{eq:dty} define the ODE system
$d\state/dz = \mathbf{F}(\state, z; \Bfield)$ that underpins both the
RK4 data generation and the physics-informed loss terms.

% ============================================================================
\section{Model Specifications}
\label{sec:models}

All architectures inherit from a common base class that provides
$z$-score input normalisation
$\tilde{\mathbf{x}} = (\mathbf{x} - \bm{\mu}_x) / \bm{\sigma}_x$
and output denormalisation
$\hat{\mathbf{y}} = \bm{\sigma}_y \tilde{\mathbf{y}} + \bm{\mu}_y$,
computed from the training set.  All trainable parameters use Xavier
uniform initialisation~\cite{Glorot2010} unless otherwise stated.

\subsection{MLP (Multi-Layer Perceptron)}
\label{sec:mlp}

The MLP is a standard feedforward network:
\begin{equation}
  \hat{\mathbf{y}}
  = W_{L+1}\,
    \sigma\!\bigl(W_L \cdots \sigma(W_1 \tilde{\mathbf{x}} + b_1) \cdots + b_L\bigr)
    + b_{L+1},
  \label{eq:mlp}
\end{equation}
where $\sigma$ is the SiLU (Swish) activation~\cite{SiLU}
$\sigma(x) = x \cdot \text{sigmoid}(x)$, and output denormalisation is
applied after the final linear layer.  Dropout is set to zero for all
\texttt{gen\_2} configurations.

The training loss is purely data-driven:
\begin{equation}
  \loss_\text{MLP} = \frac{1}{4N}\sum_{i=1}^{N}\sum_{k=1}^{4}
    \left(\frac{\hat{y}_{ik} - y_{ik}}{\sigma_k}\right)^{\!2},
  \label{eq:loss_mlp}
\end{equation}
where $\sigma_k$ is the standard deviation of output component~$k$ in
the training set ($k \in \{x,y,t_x,t_y\}$).

\paragraph{Component-weighted MSE.}
The network output is denormalised before loss computation.  Because the
physical magnitudes of the four output components span many orders of
magnitude---$|\Delta x| \sim \mathcal{O}(1)\;\mm$, $|\Delta t_x| \sim
\mathcal{O}(10^{-3})$---a plain MSE in physical space assigns the
position components $\sim10^6$ times more gradient weight than the slopes.
Dividing each component error by its training-set standard deviation
$\sigma_k$ is equivalent to computing MSE in normalised (z-score) space,
which restores gradient balance across all four components.  This is
consistent with the non-dimensionalisation approach recommended for
multi-scale PINN losses~\cite{Wang2021,Krishnapriyan2021}.

\subsection{PINN (Physics-Informed Neural Network)}
\label{sec:pinn}

The PINN uses a residual formulation that guarantees satisfaction of the
initial condition (IC) by construction.  The architecture decomposes the
prediction into a physics-based baseline (straight-line extrapolation)
plus a learned correction modulated by the fractional propagation
distance $\zeta = z/\dz \in [0, 1]$:
\begin{equation}
  \hat{\mathbf{y}}(\zeta) = \mathbf{y}_\text{base}(\zeta)
    + \zeta \cdot \bm{\delta}_{\btheta}(\mathbf{x}_0),
  \label{eq:pinn_residual}
\end{equation}
where
\begin{equation}
  \mathbf{y}_\text{base}(\zeta) =
  \begin{pmatrix}
    x_0 + t_{x,0}\,\zeta\,\dz \\
    y_0 + t_{y,0}\,\zeta\,\dz \\
    t_{x,0} \\
    t_{y,0}
  \end{pmatrix},
  \label{eq:baseline}
\end{equation}
and $\bm{\delta}_{\btheta}$ is the output of a correction network.
At $\zeta = 0$, $\hat{\mathbf{y}} = \mathbf{y}_\text{base}(0) =
(x_0, y_0, t_{x,0}, t_{y,0})$, satisfying the IC identically.  The
correction $\bm{\delta}_{\btheta}$ is predicted by an encoder that takes
only the five-dimensional initial state
$\mathbf{x}_0 = (x_0, y_0, t_{x,0}, t_{y,0}, q/p)$ as input (excluding
$\dz$), followed by a linear correction head:
\begin{equation}
  \bm{\delta}_{\btheta} = W_\text{head}\,
    \phi_L \circ \cdots \circ \phi_1(\tilde{\mathbf{x}}_0) + b_\text{head},
  \label{eq:correction}
\end{equation}
where each $\phi_l(\mathbf{h}) = \tanh(W_l \mathbf{h} + b_l)$.  The
$\tanh$ activation is used instead of SiLU to ensure smooth derivatives
for automatic differentiation of the PDE residual via
\texttt{torch.autograd.grad} with \texttt{create\_graph=True}.

The PINN training loss is:
\begin{equation}
  \loss_\text{PINN} = \loss_\text{data}
    + \lambda_\text{ic}\,\loss_\text{ic}
    + \lambda_\text{pde}\,\loss_\text{pde},
  \label{eq:loss_pinn}
\end{equation}
where each term is non-dimensionalised by the output standard deviations
$\sigma_k$ from the training set, following the gradient-balancing
prescription of Wang et al.~\cite{Wang2021}:
\begin{itemize}
  \item $\loss_\text{data}$ is the component-weighted endpoint MSE
    (Eq.~\ref{eq:loss_mlp}), identical across all three model types;

  \item $\loss_\text{ic}$ is the initial-condition loss, computed in the
    same normalised space:
    \begin{equation}
      \loss_\text{ic} = \frac{1}{4N}\sum_{i=1}^{N}\sum_{k=1}^{4}
        \left(\frac{\hat{y}_{ik}(\zeta{=}0) - y_{0,ik}}{\sigma_k}
        \right)^{\!2}.
      \label{eq:loss_ic}
    \end{equation}
    Because of the residual formulation (Eq.~\ref{eq:pinn_residual}),
    $\hat{\mathbf{y}}(\zeta{=}0) \equiv \mathbf{y}_0$ exactly, so
    $\loss_\text{ic} = 0$ by construction and serves as a sanity check.

  \item $\loss_\text{pde}$ is the PDE residual loss.  Define the raw
    residual vector at collocation point~$k$ and sample~$i$:
    \begin{equation}
      \mathbf{r}_{ik} = \frac{\partial\hat{\mathbf{y}}}{\partial z}
        \bigg|_{\zeta_k}
        - \mathbf{F}\bigl(\hat{\mathbf{y}}(\zeta_k),\;
          z_0 + \zeta_k\dz;\; \Bfield\bigr),
      \label{eq:pde_raw_residual}
    \end{equation}
    where $\mathbf{F}$ is the right-hand side of the ODE system
    (Eqs.~\ref{eq:dx}--\ref{eq:dty}).  The raw residuals for each
    component have different physical dimensions and magnitudes:
    $r^{(x)} \sim \mathcal{O}(t_x) \sim 0.1$ (dimensionless),
    while $r^{(t_x)} \sim \mathcal{O}(\kappa B) \sim 10^{-5}\;\mm^{-1}$.
    To bring all four components to $\mathcal{O}(1)$, we
    non-dimensionalise by the characteristic derivative scale
    $\sigma_k / \dz_i$:
    \begin{equation}
      \bar{r}_{ik,j} = \frac{r_{ik,j}}{\sigma_j / \dz_i},
      \qquad j \in \{x, y, t_x, t_y\},
      \label{eq:pde_nondim}
    \end{equation}
    which follows from introducing $\bar{s}_j \equiv s_j / \sigma_j$,
    $\bar{z} \equiv z / \dz$, so that $d\bar{s}_j/d\bar{z} =
    (\dz/\sigma_j)(ds_j/dz)$ is $\mathcal{O}(1)$ when the trajectory is
    well-described by the training distribution.  The PDE loss is then:
    \begin{equation}
      \loss_\text{pde} = \frac{1}{4\,n_c\,N}\sum_{k=1}^{n_c}\sum_{i=1}^{N}
        \sum_{j=1}^{4} \bar{r}_{ik,j}^{\,2},
      \label{eq:loss_pde}
    \end{equation}
    evaluated at $n_c = 8$ collocation points uniformly spaced in
    $\zeta \in [0.1, 1.0]$.  Because $\dz_i$ varies per sample, the
    non-dimensionalisation automatically applies looser absolute
    tolerances for long propagation paths while maintaining the same
    relative tolerance.
\end{itemize}

\paragraph{Note on PDE residual composition.}
All four ODE equations are enforced: two kinematic ($dx/dz = t_x$,
$dy/dz = t_y$) and two Lorentz
(Eqs.~\ref{eq:dtx}--\ref{eq:dty}).  Although the kinematic residuals
provide redundant gradient signal on the slopes (which are already
trained by $\loss_\text{data}$), they enforce internal consistency of
the trajectory---i.e.\ that the network's position and slope outputs
are mutually consistent along the full path, not just at the endpoints.
This is standard practice in the PINN literature~\cite{Raissi2019}.

\paragraph{Note on clamping.}
No output clamping is applied before computing the PDE residual.
Earlier code used \texttt{torch.clamp} to restrict predictions to a
physical range, but this breaks the \texttt{autograd} computation graph
(zero gradient outside the clamp bounds), silently zeroing the PDE loss
for out-of-range predictions.

At each collocation point, the PDE residual requires: one forward pass to
obtain $\hat{\mathbf{y}}(\zeta_k)$, four calls to
\texttt{torch.autograd.grad} to obtain $\partial\hat{y}_j/\partial\zeta$
for $j = 1, \ldots, 4$, a conversion
$\partial\hat{\mathbf{y}}/\partial z = (\partial\hat{\mathbf{y}}/\partial\zeta) / \dz$,
a magnetic field evaluation at
$(x_k, y_k, z_0 + \zeta_k \dz)$, and the Lorentz residual from
Eqs.~\eqref{eq:dx}--\eqref{eq:dty}.

The correction head weights are initialised uniformly in $[-0.01, 0.01]$
with zero bias, so that the network begins at the straight-line baseline.
A learnable scalar \texttt{field\_scale} is included but initialised to~1.

\subsection{RK\_PINN (Runge-Kutta Physics-Informed Neural Network)}
\label{sec:rkpinn}

The RK\_PINN is a multi-stage architecture inspired by the RK4 numerical
integrator~\cite{Stiasny2021, Zhai2023}.  A shared backbone extracts
features from the normalised six-dimensional input; four stage-specific
heads then predict the state at intermediate positions
$\zeta_s \in \{0.25, 0.5, 0.75, 1.0\}$:
\begin{align}
  \mathbf{f} &= \text{Backbone}(\tilde{\mathbf{x}}), \label{eq:rk_backbone} \\
  \hat{\mathbf{y}}_s &= \text{Head}_s([\mathbf{f};\, \zeta_s]),
    \quad s = 1, \ldots, 4, \label{eq:rk_heads} \\
  \hat{\mathbf{y}} &= \sum_{s=1}^{4} w_s\, \hat{\mathbf{y}}_s,
    \qquad w_s = \frac{e^{\alpha_s}}{\sum_{s'} e^{\alpha_{s'}}},
    \label{eq:rk_weights}
\end{align}
where $[\cdot;\cdot]$ denotes concatenation, $\alpha_s$ are learnable
logits initialised to $\log(c_s)$ with $c_s$ from the RK4 Butcher
tableau $(c_1, c_2, c_3, c_4) = (1, 2, 2, 1)/6$, and the softmax
ensures $\sum w_s = 1$.

The backbone consists of the first $L-1$ layers of the specified
architecture; each stage head consists of a hidden layer of size
$d_{L}$ (the last element of the \texttt{hidden\_dims} list) followed
by a linear output layer to four dimensions.  For
\texttt{hidden\_dims}~$= [96, 96]$, the backbone has one layer
(Linear$(6, 96)$ + $\tanh$) and each head has
Linear$(97, 96)$ + $\tanh$ + Linear$(96, 4)$.

The training loss mirrors the PINN:
\begin{equation}
  \loss_\text{RK\_PINN} = \loss_\text{data}
    + \lambda_\text{ic}\,\loss_\text{ic}
    + \lambda_\text{pde}\,\loss_\text{pde}.
  \label{eq:loss_rkpinn}
\end{equation}
All three terms use the \emph{same non-dimensionalisation} as the PINN
(Eqs.~\ref{eq:loss_mlp}, \ref{eq:loss_ic}, \ref{eq:loss_pde}), ensuring
that $\lambda_\text{pde}$ has a consistent physical meaning across model
types.
The IC loss is computed as the normalised MSE between the initial state and the
average of all four stage-head predictions evaluated at $\zeta = 0$.
The PDE residual is computed at each of the four stage positions using
the non-dimensional scaling $\sigma_k / \dz_i$ (Eq.~\ref{eq:pde_nondim}).
Critically, the backbone features are \emph{detached} before being
passed to the stage heads for the PDE gradient computation---gradients
propagate only through the stage heads, preventing interference between
the data loss (which trains the full network) and the PDE loss (which
trains only the heads).

\subsection{Parameter Counts and Memory Footprint}
\label{sec:param_counts}

Table~\ref{tab:param_counts} reports the trainable parameter counts
computed directly from the architecture code.  All parameters are stored
in single-precision (float32, \SI{4}{\byte} each), and the memory
footprint is $4 \times N_\text{params}$~bytes.

\begin{table}[htbp]
\centering
\caption{Trainable parameter counts and memory footprint for each
model configuration, computed from the architecture source code.  The
CUDA \texttt{\_\_constant\_\_} memory limit is \SI{64}{\kilo\byte}
(\num{16384} float32 parameters).}
\label{tab:param_counts}
\begin{tabular}{llrrr}
Config & Type & Dims & Params & KB \\
\hline
\texttt{mlp/tiny}   & MLP & $[64,64]$   & \num{4868}  & 19.0 \\
\texttt{mlp/small}  & MLP & $[96,96]$   & \num{10372} & 40.5 \\
\texttt{mlp/medium} & MLP & $[112,112]$ & \num{13892} & 54.3 \\
\texttt{pinn/*}     & PINN & $[96,96]$  & \num{10277}\textsuperscript{a}
                                                       & 40.1 \\
\texttt{rk\_pinn/*} & RK\_PINN & $[96,96]$
                     & \num{39860}\textsuperscript{b} & 155.7 \\
\end{tabular}
\begin{flushleft}
\textsuperscript{a} Encoder input is 5-dimensional (no $\dz$); includes
one learnable \texttt{field\_scale} parameter. \\
\textsuperscript{b} Four separate stage heads each contain
Linear$(97, 96)$ + Linear$(96, 4)$; total is $672 + 4 \times 9796 + 4 =
\num{39860}$.  This \emph{exceeds} the \SI{64}{\kilo\byte}
\texttt{\_\_constant\_\_} memory limit and is therefore relevant only to
RQ2 (scientific interest), not RQ1 (engineering deployment).
\end{flushleft}
\end{table}

\paragraph{Remark.}
The PINN encoder takes a 5-dimensional input (initial state without
$\dz$), yielding 96 fewer first-layer weights than the MLP; it also
includes one additional scalar parameter (\texttt{field\_scale}).  The
RK\_PINN, by design, replicates the output head four times, resulting in
approximately four times the parameter count of a single-head model with
the same hidden dimensions.  This means the ``same architecture size''
comparison between MLP, PINN, and RK\_PINN at \texttt{hidden\_dims}~$=
[96, 96]$ is \emph{not} a controlled comparison of parameter count.
Section~\ref{sec:confounds} discusses this confound in detail.

% ============================================================================
\section{Experimental Design}
\label{sec:experimental_design}

\subsection{Configuration Matrix}

Table~\ref{tab:configs} specifies all nine training runs.  Each
configuration is defined by a JSON file in
\texttt{experiments/gen\_2/configs/}.

\begin{table}[htbp]
\centering
\caption{Complete configuration matrix for the \texttt{gen\_2} experiment.
All runs use seed~42, weight decay $10^{-4}$, maximum 300 epochs,
gradient clipping at 1.0, and cosine annealing with 5-epoch linear warmup.
Physics loss weights $\lambda_\text{pde}$ and $\lambda_\text{ic}$ apply
only to PINN and RK\_PINN models.}
\label{tab:configs}
\footnotesize
\begin{tabular}{llccccccccc}
Config & Type & Hidden & Act. & Batch & lr & Pat. &
  $\lambda_\text{pde}$ & $\lambda_\text{ic}$ & $n_c$ & Field \\
\hline
\texttt{mlp/tiny}       & MLP     & $[64,64]$   & SiLU & 4096 & $10^{-3}$
  & 30 & --- & --- & --- & --- \\
\texttt{mlp/small}      & MLP     & $[96,96]$   & SiLU & 4096 & $10^{-3}$
  & 30 & --- & --- & --- & --- \\
\texttt{mlp/medium}     & MLP     & $[112,112]$ & SiLU & 4096 & $10^{-3}$
  & 30 & --- & --- & --- & --- \\
\texttt{pinn/lam0.1}    & PINN    & $[96,96]$   & tanh & 1024 & $5 \times 10^{-4}$
  & 40 & 0.1  & 0.1 & 8 & interp. \\
\texttt{pinn/lam1.0}    & PINN    & $[96,96]$   & tanh & 1024 & $5 \times 10^{-4}$
  & 40 & 1.0  & 0.1 & 8 & interp. \\
\texttt{pinn/lam10.0}   & PINN    & $[96,96]$   & tanh & 1024 & $5 \times 10^{-4}$
  & 40 & 10.0 & 0.1 & 8 & interp. \\
\texttt{rk\_pinn/lam0.1}  & RK\_PINN & $[96,96]$ & tanh & 1024 & $5 \times 10^{-4}$
  & 40 & 0.1  & 1.0 & 4\textsuperscript{$\dagger$} & interp. \\
\texttt{rk\_pinn/lam1.0}  & RK\_PINN & $[96,96]$ & tanh & 1024 & $5 \times 10^{-4}$
  & 40 & 1.0  & 1.0 & 4\textsuperscript{$\dagger$} & interp. \\
\texttt{rk\_pinn/lam10.0} & RK\_PINN & $[96,96]$ & tanh & 1024 & $5 \times 10^{-4}$
  & 40 & 10.0 & 1.0 & 4\textsuperscript{$\dagger$} & interp. \\
\end{tabular}
\begin{flushleft}
\textsuperscript{$\dagger$} RK\_PINN evaluates the PDE residual at the
four fixed stage positions $\zeta \in \{0.25, 0.5, 0.75, 1.0\}$ rather
than at $n_c$ uniformly spaced collocation points.
\end{flushleft}
\end{table}

\subsection{Comparisons and Controls}

The nine configurations enable the following structured comparisons:

\paragraph{MLP size scaling (RQ1).}
\texttt{mlp/tiny} $\to$ \texttt{mlp/small} $\to$ \texttt{mlp/medium}
tests how accuracy improves with model capacity within the
\SI{64}{\kilo\byte} budget.  The controlled variables are activation
(SiLU), batch size (4096), and learning rate ($10^{-3}$).

\paragraph{MLP vs.\ PINN at matched architecture (RQ2).}
\texttt{mlp/small}~$[96,96]$ vs.\ \texttt{pinn/lam*}~$[96,96]$ compares
data-driven and physics-informed training.  The hidden dimensions match,
but several variables are \emph{not} controlled (see
Sec.~\ref{sec:confounds}).

\paragraph{MLP vs.\ RK\_PINN at matched hidden dimensions (RQ2).}
\texttt{mlp/small}~$[96,96]$ vs.\ \texttt{rk\_pinn/lam*}~$[96,96]$.
The hidden dimensions match but the RK\_PINN has $\approx 4\times$ more
parameters due to four stage heads (Table~\ref{tab:param_counts}).

\paragraph{PINN $\lambda_\text{pde}$ sweep (RQ2).}
$\lambda_\text{pde} \in \{0.1, 1.0, 10.0\}$ at fixed
$\lambda_\text{ic} = 0.1$ tests the sensitivity of the PINN to the
physics loss weight.

\paragraph{RK\_PINN $\lambda_\text{pde}$ sweep (RQ2).}
$\lambda_\text{pde} \in \{0.1, 1.0, 10.0\}$ at fixed
$\lambda_\text{ic} = 1.0$ for the RK\_PINN architecture.

\paragraph{PINN vs.\ RK\_PINN (RQ2).}
At each $\lambda_\text{pde}$ value, the PINN and RK\_PINN can be
compared to assess whether the multi-stage RK structure provides
additional benefit.  This comparison is confounded by the parameter
count difference.

% ============================================================================
\section{Training Protocol}
\label{sec:training}

\subsection{Optimiser and Schedule}

All models are trained with AdamW~\cite{AdamW}
($\beta_1 = 0.9$, $\beta_2 = 0.999$, weight decay $= 10^{-4}$).  The
learning rate follows a cosine annealing schedule with a 5-epoch linear
warmup:
\begin{equation}
  \text{lr}(t) =
  \begin{cases}
    \text{lr}_0 \cdot t / 5, & t \leq 5, \\[4pt]
    \dfrac{\text{lr}_0}{2}
    \Bigl(1 + \cos\!\bigl(\pi \cdot \tfrac{t - 5}{T - 5}\bigr)\Bigr),
    & t > 5,
  \end{cases}
  \label{eq:schedule}
\end{equation}
where $T = 300$ is the maximum number of epochs.  Gradient norms are
clipped to 1.0.

\subsection{Early Stopping}

Training terminates early if the validation loss does not improve by more
than $\delta_\text{min} = 10^{-7}$ for a number of consecutive epochs
specified by the patience parameter: 30 for MLPs, 40 for PINN and
RK\_PINN.  The extended patience for physics-informed models accounts for
the slower convergence typically observed when balancing multiple loss
terms.

\subsection{Checkpointing}

Model weights are saved every 25 epochs and at the best validation loss.
The best-epoch checkpoint is used for all evaluation.

\subsection{Infrastructure}

Training is executed on the Nikhef HTCondor cluster
(\texttt{taai-007.nikhef.nl}) via the submission file
\texttt{condor/submit.sub}, which dispatches one GPU job per
configuration.  GPU nodes provide either NVIDIA Tesla V100-PCIE-32GB or
NVIDIA L40S accelerators.  The \texttt{+JobCategory = "long"} policy
allows extended wall-clock time.  The software environment is:
\begin{itemize}
  \item Conda environment \texttt{TE} (Python~3.10),
  \item PyTorch 2.9.1 with CUDA 12.8,
  \item NumPy 2.2.6, MLflow 3.10.1.
\end{itemize}
Full dependency pinning is recorded in \texttt{environment.yml} and
\texttt{requirements.txt}.

\subsection{Experiment Tracking}

\begin{sloppypar}
All runs are logged to MLflow under the experiment name
\nolinkurl{gen_2_track_extrapolation}.  Logged artefacts include:
training and validation loss per epoch (total and per-component for
PINN/RK\_PINN), learning rate schedule, model configuration,
best-epoch checkpoint, and test-set evaluation metrics.  TensorBoard
summaries are written in parallel.
\end{sloppypar}

% ============================================================================
\section{Evaluation Protocol}
\label{sec:evaluation}

All metrics are computed on the held-out test set (\num{5e6} samples,
disjoint from training and validation data).

\subsection{Primary Metrics}

\paragraph{Position error.}
\begin{equation}
  \|\Delta \mathbf{r}\| = \sqrt{(\hat{x}_f - x_f)^2 + (\hat{y}_f - y_f)^2},
  \label{eq:pos_err}
\end{equation}
reported as mean, standard deviation, 95th percentile, and 99th percentile.

\paragraph{Slope error.}
\begin{equation}
  \|\Delta \mathbf{s}\| = \sqrt{(\hat{t}_{x,f} - t_{x,f})^2
    + (\hat{t}_{y,f} - t_{y,f})^2},
  \label{eq:slope_err}
\end{equation}
reported as mean and 95th percentile.

\paragraph{Per-component errors.}
$|\Delta x|$, $|\Delta y|$, $|\Delta t_x|$, $|\Delta t_y|$ reported
as means.

\subsection{Secondary Metrics}

\paragraph{Jacobian (transport matrix).}
The $4 \times 5$ Jacobian
\begin{equation}
  J_{ij} = \frac{\partial \hat{y}_i}{\partial x_{0,j}},
  \quad i \in \{x, y, t_x, t_y\},\;
  j \in \{x, y, t_x, t_y, q/p\},
  \label{eq:jacobian}
\end{equation}
is computed via \texttt{torch.autograd.grad} on 1000 randomly selected
test samples.  The Frobenius norm $\|J\|_F$ is reported as mean and
standard deviation.

\paragraph{Remark on Jacobian discriminative power.}
In \texttt{gen\_1}, all models produced $\|J\|_F \approx 5650 \pm 4200$
with minimal inter-model variation.  This suggests the metric is
dominated by the large-$\dz$ regime where the Jacobian magnitude is
inherently large.  A $\dz$-binned analysis (Sec.~\ref{sec:analysis_plan})
may be more informative.

\subsection{Success Criteria}
\label{sec:success}

\paragraph{RQ1.}
A model is considered viable for Allen HLT1 deployment if:
\begin{enumerate}
  \item it fits within \SI{64}{\kilo\byte} (\num{16384} float32 parameters),
  \item its mean position error $\langle\|\Delta\mathbf{r}\|\rangle
    < \SI{0.5}{\mm}$ on the full test set,
  \item its 95th-percentile position error $< \SI{2}{\mm}$,
  \item its mean slope error $\langle\|\Delta\mathbf{s}\|\rangle < 0.005$.
\end{enumerate}
These thresholds are informed by \texttt{gen\_1} results and by the
requirement that the Kalman filter converge to the same track parameters
as with the RK4 extrapolator.  A full integration test in Allen is beyond
the scope of this study but is the logical next step if these criteria
are met.

\paragraph{RQ2.}
Physics-informed training is considered beneficial if, at the same
\texttt{hidden\_dims}~$= [96, 96]$:
\begin{enumerate}
  \item the best PINN or RK\_PINN configuration achieves lower mean
    position error than \texttt{mlp/small}, after accounting for the
    confounds listed in Sec.~\ref{sec:confounds}; or
  \item the physics-informed model shows qualitatively better
    generalisation in the low-statistics tails of the parameter space
    (low $p$, large $\dz$, or extreme slopes).
\end{enumerate}

% ============================================================================
\section{Analysis Plan}
\label{sec:analysis_plan}

\subsection{Aggregate Comparisons}

\begin{enumerate}
  \item \textbf{Summary table}: All nine models ranked by mean position
    error, 95th-percentile position error, mean slope error, and
    best epoch.  Include training wall-clock time per epoch.
  \item \textbf{MLP scaling plot}: Mean and 95th-percentile position
    error vs.\ parameter count for the three MLP sizes.
  \item \textbf{$\lambda_\text{pde}$ sensitivity}: Position error vs.\
    $\lambda_\text{pde}$ for PINN and RK\_PINN families.
\end{enumerate}

\subsection{$\dz$-Binned Analysis}

To assess performance across the full propagation range, the test set will
be binned in $\dz$:
\begin{itemize}
  \item $[25, 100]$~mm (VELO inter-sensor),
  \item $[100, 500]$~mm (short-range),
  \item $[500, 2000]$~mm (medium-range),
  \item $[2000, 5000]$~mm (long-range),
  \item $[5000, 10000]$~mm (full-detector).
\end{itemize}
For each bin: mean, 95th percentile, and 99th percentile of
$\|\Delta\mathbf{r}\|$ and $\|\Delta\mathbf{s}\|$.  This is essential to
verify that the extended $\dz$ coverage (\texttt{gen\_2}) does not degrade
short-$\dz$ performance relative to \texttt{gen\_1}.

\subsection{Momentum-Binned Analysis}

Low-momentum tracks ($p < \SI{5}{\giga\eV}/c$) experience the strongest
magnetic deflection and are most challenging for the network.  The test
set will be binned in $p$:
\begin{itemize}
  \item $[1, 5]$~GeV/$c$ (strongly bent),
  \item $[5, 20]$~GeV/$c$ (moderately bent),
  \item $[20, 200]$~GeV/$c$ (weakly bent).
\end{itemize}

\subsection{Diagnostic Plots}
\label{sec:diagnostic_plots}

The following visualisations will be produced for each model:
\begin{enumerate}
  \item Training and validation loss curves (total, and per-component for
    PINN/RK\_PINN).
  \item 2D heatmap of $\|\Delta\mathbf{r}\|$ in the $(\dz, p)$ plane.
  \item Scatter plot of predicted vs.\ true $x_f$ and $t_{x,f}$ for a
    random subsample.
  \item Residual distributions $\Delta x$, $\Delta y$, $\Delta t_x$,
    $\Delta t_y$ (histograms with Gaussian fits).
  \item Jacobian Frobenius norm vs.\ $\dz$ (binned).
  \item RK\_PINN stage weights evolution during training
    (initialised at $[1,2,2,1]/6$; do they remain near RK4?).
  \item PINN IC loss over training (should remain near zero if residual
    formulation works correctly).
\end{enumerate}

\subsection{Gen\_1 vs.\ Gen\_2 Comparison}

The \texttt{mlp/tiny} configuration $[64, 64]$ is shared between
\texttt{gen\_1} and \texttt{gen\_2} (though trained on different data).
Direct comparison of position and slope errors will quantify the impact
of extending $\dz$ coverage.  The \texttt{gen\_1} baseline values are:
\begin{center}
\begin{tabular}{lcc}
\toprule
Metric & gen\_1 & gen\_2 \\
\midrule
$\langle\|\Delta\mathbf{r}\|\rangle$ (mm) & 0.491 & (this study) \\
$\|\Delta\mathbf{r}\|_{95\%}$ (mm) & 1.452 & (this study) \\
$\langle\|\Delta\mathbf{s}\|\rangle$ & 0.00491 & (this study) \\
Best epoch & 105 & (this study) \\
\bottomrule
\end{tabular}
\end{center}

% ============================================================================
\section{Known Limitations and Confounds}
\label{sec:confounds}

This section catalogues the known limitations of the experimental design.
These must be considered when interpreting results.

\begin{enumerate}
  \item \textbf{Activation function mismatch.}
    MLPs use SiLU; PINN and RK\_PINN use $\tanh$.  The $\tanh$ choice is
    motivated by the requirement for smooth second-order derivatives in
    the PDE loss (SiLU's second derivative has a discontinuity at
    $x \to -\infty$ in the limit).  However, SiLU generally outperforms
    $\tanh$ for regression tasks due to its non-saturating positive
    branch~\cite{SiLU}.  If a physics-informed model outperforms the MLP,
    the improvement cannot be unambiguously attributed to the physics loss;
    conversely, if the MLP outperforms, the physics loss may be beneficial
    but masked by the inferior activation.  A controlled comparison would
    require running the MLP with $\tanh$ activation.

  \item \textbf{Batch size mismatch.}
    MLPs train with batch size 4096; PINN and RK\_PINN use 1024 due to
    the memory overhead of the physics loss computation (8 collocation
    points $\times$ forward pass $\times$ 4 autograd calls).  At a
    training set of \num{40e6} samples, one epoch comprises $\sim$9766
    steps at batch size 4096 vs.\ $\sim$39063 steps at batch size 1024.
    The physics-informed models therefore see $4\times$ more optimiser
    updates per epoch, which may independently affect convergence.

  \item \textbf{Learning rate mismatch.}
    MLPs use $\text{lr} = 10^{-3}$; PINN/RK\_PINN use $5 \times 10^{-4}$.
    The lower rate was chosen for training stability when balancing
    multiple loss terms, but it constitutes an uncontrolled variable.

  \item \textbf{Parameter count mismatch (RK\_PINN).}
    As shown in Table~\ref{tab:param_counts}, the RK\_PINN at
    \texttt{hidden\_dims}~$= [96, 96]$ has $\sim$\num{40000} parameters
    ($\sim$\SI{156}{\kilo\byte}), approximately $4\times$ the MLP and
    PINN at the same hidden dimensions.  This makes direct
    accuracy-per-parameter comparison misleading.  The RK\_PINN results
    are primarily of scientific interest (RQ2) and cannot address RQ1
    (GPU deployment) without architectural changes.

  \item \textbf{PINN training cost.}
    The PDE loss with 8 collocation points requires $\sim$8 additional
    forward passes and 32 autograd calls per batch.  PINN epochs take
    $\sim$\SI{50}{\minute} vs.\ $\sim$\SI{7}{\minute} for the MLP on a
    V100.  Per-epoch training cost differs by approximately $7\text{--}8\times$.

  \item \textbf{No explicit RK4 timing baseline.}
    The ground truth is the RK4 output, so the ``baseline error'' is zero
    by definition.  The relevant question is whether the NN error is small
    enough for the downstream Kalman filter, which requires an integration
    test in Allen that is beyond the scope of this protocol.

  \item \textbf{Single polarity.}
    All data use MagDown ($q = -1$).  MagUp data are not included.  The
    field map symmetry $B_y(x, y, z; \text{MagUp}) = -B_y(x, y, z;
    \text{MagDown})$ implies that a model trained on MagDown should work
    for MagUp with a sign flip of $q/p$, but this is not tested.

  \item \textbf{No material interactions.}
    The RK4 integrator conserves momentum ($dq/p\,/\,dz = 0$).  Real
    tracks lose energy and scatter in detector material.  The trained
    models learn vacuum propagation only.

  \item \textbf{Jacobian metric sensitivity.}
    The Jacobian Frobenius norm $\|J\|_F$ showed poor discriminative
    power in \texttt{gen\_1}, with all models yielding similar values.
    This may be because $\|J\|_F$ is dominated by the $\partial x_f /
    \partial x_0 \approx 1$ and $\partial y_f / \partial y_0 \approx 1$
    diagonal elements at small $\dz$, while the physically interesting
    off-diagonal elements (e.g., $\partial x_f / \partial (q/p)$) are
    comparatively small.
\end{enumerate}

% ============================================================================
\section{Reproducibility}
\label{sec:reproducibility}

\subsection{Random Seeds}

\begin{sloppypar}
A single global seed (42) controls all sources of randomness:
\nolinkurl{random.seed}, \nolinkurl{numpy.random.seed},
\nolinkurl{torch.manual_seed}, and
\nolinkurl{torch.cuda.manual_seed_all}.
Additionally,
\nolinkurl{torch.backends.cudnn.deterministic} is set to \texttt{True} and
\nolinkurl{torch.backends.cudnn.benchmark} to \texttt{False}.
Data generation uses per-batch seeds as specified in
Eq.~\eqref{eq:seed}.
\end{sloppypar}

\subsection{File Manifest}

All files required to reproduce this experiment are contained within
\texttt{experiments/gen\_2/}:
\begin{center}
\begin{tabular}{ll}
\toprule
Path & Contents \\
\midrule
\texttt{configs/mlp/*.json}     & MLP configurations (3 files) \\
\texttt{configs/pinn/*.json}    & PINN configurations (3 files, legacy) \\
\texttt{configs/rk\_pinn/*.json}& RK\_PINN configurations (3 files, legacy) \\
\texttt{configs/v2\_fixes/pinn\_v2/*.json}    & PINN\_v2 configs (3 files, 2026-04-21) \\
\texttt{configs/v2\_fixes/neural\_rk4/*.json} & NeuralRK4 configs (3 files, 2026-04-21) \\
\texttt{models/train.py}        & Training script (Fix F added 2026-04-21) \\
\texttt{models/architectures.py}& Model definitions incl.\ \texttt{PINN\_v2}, \texttt{NeuralRK4} \\
\texttt{utils/rk4\_propagator.py}& RK4 integrator \\
\texttt{utils/magnetic\_field.py}& Unified field module \\
\texttt{data/merge\_batches.py} & Batch merger \\
\texttt{condor/submit\_datagen.sub} & Data generation jobs \\
\texttt{condor/run\_datagen.sh}     & Data generation worker \\
\texttt{condor/submit.sub}          & Original training jobs (9 entries) \\
\texttt{condor/submit\_v2\_fixes.sub} & v2 fixes training jobs (6 entries) \\
\texttt{condor/run\_training.sh}    & Training worker \\
\texttt{condor/jobs.txt}            & Original job list \\
\texttt{condor/jobs\_v2\_fixes.txt} & v2 fixes job list \\
\texttt{ISSUES\_WITH\_PINNS\_AND\_RK\_PINNS.md} & Audit + resolution report (2026-04-21) \\
\bottomrule
\end{tabular}
\end{center}

\subsection{Software Environment}

The full environment specification is frozen in:
\begin{itemize}
  \item \texttt{environment.yml} --- Conda environment with all packages
    pinned to exact versions,
  \item \texttt{requirements.txt} --- pip-installable subset with version
    pins.
\end{itemize}
Key versions: Python~3.10, PyTorch~2.9.1+cu128, NumPy~2.2.6,
MLflow~3.10.1, SciPy~1.15.3.

\subsection{MLflow Tracking}

\begin{sloppypar}
All training runs are logged to the local MLflow tracking server
under experiment \texttt{gen\_2\_track\_\allowbreak extrapolation}.  The MLflow
database is stored at \texttt{experiments/\allowbreak gen\_2/\allowbreak mlruns/}.
Run artefacts include the best model checkpoint, training
configuration, per-epoch metrics, and test-set evaluation results.
\end{sloppypar}

% ============================================================================
\section{Architectural Fixes: \texttt{PINN\_v2} and \texttt{NeuralRK4}
(2026-04-21 Revision)}
\label{sec:v2_fixes}

After the initial \texttt{gen\_2} submission completed, a theoretical audit
(companion document
\texttt{experiments/gen\_2/ISSUES\_WITH\_PINNS\_AND\_RK\_PINNS.md}) identified
three architectural pathologies that prevented the legacy PINN and RK\_PINN
models (Secs.~\ref{sec:pinn}, \ref{sec:rkpinn}) from converging to a useful
operating point on their own terms. This section documents the fixes applied
on 2026-04-21 and the two replacement architectures --- \texttt{PINN\_v2}
and \texttt{NeuralRK4} --- that now supersede the original PINN and RK\_PINN
in the gen\_2 submission.

\subsection{Diagnosed Pathologies}
\label{sec:v2_pathologies}

\begin{enumerate}
  \item \textbf{PINN is structurally linear in $\zeta = z / \dz$.}
    The legacy PINN envelope (Eq.~\ref{eq:pinn_residual}),
    $\hat{\state}(\zeta) = \state_\text{base}(\zeta) + \zeta\,\bm{\delta}_{\btheta}(\state_0)$,
    contains no $\zeta$-dependent nonlinearity because the encoder is fed
    only the initial state $\state_0$.  The trajectory class is therefore
    affine in $\zeta$; since the Lorentz ODE is nonlinear in $\state$ and
    nonlinear in $z$ through $\Bfield(x,y,z)$, the PDE residual
    (Eq.~\ref{eq:pde_raw_residual}) has a \emph{structural lower bound}
    that the optimiser cannot reduce.  Empirically, the PINN
    $\lambda_\text{pde}$ response was non-monotonic.
  \item \textbf{PDE residual cost was $N_c \times 4$ reverse-mode backwards
    per batch.}  With $N_c = 8$ collocation points and four output
    components, the legacy implementation performed 32
    \texttt{torch.autograd.grad} calls per batch.  Because the input
    $\zeta$ is one-dimensional and the output four-dimensional, the
    Jacobian is $4 \times 1$ and forward-mode automatic differentiation
    (JVP) recovers all four components in a single sweep.
  \item \textbf{RK\_PINN was not a Runge--Kutta method.}
    The softmax-weighted combination (Eq.~\ref{eq:rk_weights}) averages
    \emph{states} evaluated at $\zeta \in \{0.25, 0.5, 0.75, 1.0\}$, not
    \emph{derivatives}.  Classical RK4 averages derivatives of a single
    propagating state with fixed Butcher weights $(1, 2, 2, 1)/6$ to
    produce one endpoint state; a convex combination of states at
    different $\zeta$ has no approximation-theoretic justification.
    Combined with the \texttt{features.detach()} call in the physics
    branch, this produced slope errors 43--587$\times$ worse than the MLP.
\end{enumerate}

The proposed fixes (labels follow the audit document) are:
{\small
\begin{center}
\begin{tabular}{cll}
Label & Change & Priority \\
\hline
A  & Replace \texttt{autograd.grad} loop with \texttt{torch.func.jvp} & P0 \\
B  & Stochastic collocation: $n_c$ MC draws $\zeta \sim \mathcal{U}(0,1)$ per sample & P0 \\
C  & Give PINN a $\zeta$-dependent trunk & P0 \\
D' & Rewrite RK-style model as neural ODE with RK4 integrator over exact Lorentz RHS & P1 \\
E  & Remove \texttt{features.detach()} in \texttt{RK\_PINN.compute\_physics\_loss} & P0 \\
F  & Linear warm-up $\lambda_{\text{pde},\text{ic}}(t) = \lambda^\star \min(1, (t{+}1)/5)$ & P1 \\
\end{tabular}
\end{center}
}

\subsection{\texttt{PINN\_v2} --- Fixes A + B + C + F}
\label{sec:pinn_v2}

The corrected envelope is
\begin{equation}
  \hat{\state}_\btheta(\zeta) \;=\;
  \state_\text{base}(\zeta)
  \;+\; \zeta \cdot
  \mathrm{Head}\!\Bigl(
     \mathrm{Enc}\bigl( [\tilde{\state}_0;\,\zeta] \bigr)
  \Bigr),
  \label{eq:pinn_v2_envelope}
\end{equation}
where $\state_\text{base}$ is still the straight-line baseline
(Eq.~\ref{eq:baseline}) and
$[\tilde{\state}_0;\,\zeta] \in \mathbb{R}^{6}$ replaces the
$\mathbb{R}^{5}$ input of the legacy encoder (Eq.~\ref{eq:correction}).
Because the $\zeta$-prefactor still multiplies the entire correction,
$\hat{\state}_\btheta(0) \equiv \state_0$ exactly and the IC loss
(Eq.~\ref{eq:loss_ic}) remains identically zero by construction.
The trajectory is now a smooth nonlinear function of $\zeta$, and the
structural floor on $\loss_\text{pde}$ is eliminated.

The PDE residual (Eq.~\ref{eq:pde_raw_residual}) is evaluated with a single
forward-mode JVP per Monte-Carlo collocation draw:
\begin{equation}
  \bigl(\hat{\state}(\zeta_i),\;
  \tfrac{\partial \hat{\state}}{\partial \zeta}\bigr\rvert_{\zeta_i}\bigr)
  \;=\;
  \texttt{torch.func.jvp}\bigl(
     \mathrm{forward\_at\_z},\;
     (\zeta_i,),\; (1,)
  \bigr),
  \qquad \zeta_i \sim \mathcal{U}(0,1),
  \label{eq:pinn_v2_jvp}
\end{equation}
with $n_c$ independent draws per training sample (default $n_c = 2$).
This replaces $N_c \times 4$ reverse-mode backwards with a single JVP
sweep per draw; CPU microbenchmarks show a $\sim2\times$ per-iteration
speed-up even before combining with the reduced $n_c$, and a $\sim30\times$
reduction in PDE-term cost overall.  The data and IC losses retain the
non-dimensionalisation of Eqs.~\eqref{eq:loss_mlp} and \eqref{eq:loss_ic};
the PDE residual uses the same $\sigma_k / \dz$ scaling as
Eq.~\eqref{eq:pde_nondim}.  A linear warm-up (Fix F),
\begin{equation}
  \lambda_\text{pde}(t) = \lambda^\star \cdot \min\!\bigl(1,\; (t{+}1)/T_\text{warm}\bigr),
  \qquad T_\text{warm} = 5\ \text{epochs},
  \label{eq:warmup}
\end{equation}
prevents the early-epoch data/physics conflict that caused the legacy
PINN to early-stop at epoch 10--16.

\paragraph{Parameter count.}
\texttt{PINN\_v2} with \texttt{hidden\_dims}~$= [256, 256]$ has
\num{68612} parameters, dominated by the wider trunk required to
reach the PDE-residual floor of the fully nonlinear model class.

\subsection{\texttt{NeuralRK4} --- Fix D$'$}
\label{sec:neural_rk4}

\texttt{NeuralRK4} replaces the RK\_PINN state-mixture with a genuine
fourth-order Runge--Kutta integrator of the Lorentz ODE, perturbed by a
small learned correction to the right-hand side:
\begin{equation}
  \mathbf{f}_\btheta(\state, z; q/p)
  \;=\;
  \mathbf{F}(\state, z; q/p)
  \;+\;
  e^{\alpha}\,\bigl(\bm{\sigma}_y / \dz\bigr) \odot
  \mathbf{g}_\phi(\state, z; q/p),
  \label{eq:neural_rk4_rhs}
\end{equation}
where $\mathbf{F}$ is the exact Lorentz RHS
(Eqs.~\ref{eq:dx}--\ref{eq:dty}), $\mathbf{g}_\phi$ is a small MLP, and
$\alpha$ is a single learnable log-scale initialised at $\log(10^{-3})$
so that training begins at ``pure RK4 plus an infinitesimal drift.''  The
$\bm{\sigma}_y / \dz$ factor brings the learned correction into the
physical scale of the RHS.  The state update is
\begin{equation}
  \state_{i+1} \;=\; \state_i
    + \tfrac{h}{6}\bigl( \mathbf{k}_1 + 2\mathbf{k}_2 + 2\mathbf{k}_3 + \mathbf{k}_4 \bigr),
  \quad
  \mathbf{k}_m = \mathbf{f}_\btheta\bigl(\state_i^{(m)}, z_i^{(m)}; q/p\bigr),
  \quad
  h = \dz / N_\text{RK},
  \label{eq:neural_rk4_step}
\end{equation}
iterated $N_\text{RK} \in \{1, 2\}$ times with the standard RK4 nodes
$\state_i^{(m)}$.  Butcher weights are \emph{fixed} (not learned) and
multiply \emph{derivatives}, restoring the approximation-theoretic meaning
the original RK\_PINN lacked (Sec.~\ref{sec:v2_pathologies}~(3)).

Because the physics is baked into the integrator, no auxiliary PDE or IC
loss is required: the training loss is the same component-weighted
endpoint MSE as the MLP (Eq.~\ref{eq:loss_mlp}),
$\lambda_\text{pde} = \lambda_\text{ic} = 0$.

\paragraph{Parameter count.}
\texttt{NeuralRK4} with \texttt{hidden\_dims}~$= [64, 64]$ has \num{4869}
trainable parameters (matched to \texttt{mlp/tiny}); with
\texttt{hidden\_dims}~$= [128, 128]$ it has \num{17925}.  The
$N_\text{RK}$ choice does not change the parameter count, only the
integration granularity.

\subsection{Measured Outcome on \texttt{train\_50M\_dz25.npz}}
\label{sec:v2_results}

All v2 runs share the training data, split (80/10/10), seed (42), batch
size (1024 for \texttt{PINN\_v2} and \texttt{neural\_rk4/small\_2step};
2048 otherwise), base learning rate ($5 \times 10^{-4}$), and cosine
annealing schedule of Sec.~\ref{sec:training}.  The MLP baselines were
trained on the same dataset under the schedule of
Table~\ref{tab:configs}.  Table~\ref{tab:v2_results} summarises the
test-set metrics; two v2 runs are still in progress at the time of this
edit and are marked with the latest validation-log value.

\begin{table}[htbp]
\centering
\caption{Final test metrics for the \texttt{gen\_2} v2 submission
(cluster 4262459).  MLP baselines from cluster 4262446 on the same
\texttt{train\_50M\_dz25.npz} dataset included for comparison.  All
quantities computed on the 5M-sample test split.  Starred entries are
the latest value parsed from the training log for the two runs still in
progress.}
\label{tab:v2_results}
\footnotesize
\begin{tabular}{llrrrrrr}
Family & Config & Hidden & Params & Pos [mm] & Slope [rad] &
  best ep / total & Wall [h] \\
\hline
MLP       & tiny   & $[64, 64]$           & \num{4868}  & 0.3868 & $8.5 \times 10^{-5}$  & 62 / 64 & 3.69 \\
MLP       & small  & $[128, 128]$         & \num{18692} & 0.5092 & $1.21 \times 10^{-4}$ & 16 / 46 & 2.57 \\
MLP       & medium & $[256, 256, 128]$    & \num{102788}& 0.3046 & $9.5 \times 10^{-5}$  & 20 / 40 & 1.63 \\
MLP       & large  & $[512, 512, 256]$    & \num{401412}& 0.3798 & $7.6 \times 10^{-5}$  & 14 / 44 & 1.81 \\
MLP       & wide   & $[512, 512, 256, 128]$ & \num{434180}& 0.3139 & $7.9 \times 10^{-5}$  & 24 / 42 & 1.87 \\
\hline
PINN\_v2   & $\lambda = 0.1$  & $[256, 256]$ & \num{68612} & \textbf{0.2352} & $6.8 \times 10^{-5}$  & 4 / 41  & 7.84 \\
PINN\_v2   & $\lambda = 1.0$  & $[256, 256]$ & \num{68612} & 0.3218          & $8.6 \times 10^{-5}$  & 1 / 41  & 7.87 \\
PINN\_v2   & $\lambda = 10.0$ & $[256, 256]$ & \num{68612} & $\sim 0.40$\textsuperscript{$\star$} &
                $1.5 \times 10^{-4}$\textsuperscript{$\star$} & --- & --- \\
NeuralRK4 & tiny, 1 step  & $[64, 64]$    & \num{4869}  & 0.1339 & $5.0 \times 10^{-5}$  & 40 / 42 & 3.62 \\
NeuralRK4 & small, 1 step & $[128, 128]$  & \num{17925} & \textbf{0.1251} & $\mathbf{5.0 \times 10^{-5}}$ & 40 / 42 & 3.64 \\
NeuralRK4 & small, 2 steps & $[128, 128]$ & \num{17925} & $\sim 0.128$\textsuperscript{$\star$} &
                $5.0 \times 10^{-5}$\textsuperscript{$\star$} & --- & --- \\
\end{tabular}
\end{table}

\paragraph{Before/after contrast.}
The pre-fix results on an earlier cluster (4217847, variable-\dz{}
dataset) gave RK\_PINN slope errors in
$[5.2 \times 10^{-2},\; 1.7 \times 10^{-1}]$~rad and a non-monotone PINN
$\lambda_\text{pde}$ response (see
\texttt{ISSUES\_WITH\_PINNS\_AND\_RK\_PINNS.md}~\S2).  On the v2
submission:
\begin{enumerate}
  \item \emph{Slope catastrophe eliminated.} NeuralRK4 achieves
    $5.0 \times 10^{-5}$~rad --- a three-orders-of-magnitude improvement
    over the legacy RK\_PINN --- because Butcher weights multiply the
    correct quantity (derivatives) and the backbone is no longer detached
    from the physics gradient.
  \item \emph{Monotone $\lambda_\text{pde}$ response recovered.}
    PINN\_v2 position degrades monotonically as
    $\lambda_\text{pde}$ increases ($0.235 \to 0.322 \to \sim 0.40$~mm),
    consistent with a genuine data/physics trade-off rather than the
    structural-floor pathology of the legacy PINN.
  \item \emph{MLPs beaten on accuracy per parameter.}  At
    \num{4869} parameters (matched to \texttt{mlp/tiny}), NeuralRK4
    reaches 0.134~mm, 2.3$\times$ better than the best MLP
    (\texttt{mlp/medium}, \num{102788} parameters, 0.305~mm) with
    $\sim$21$\times$ fewer parameters.  The gains at \num{17925}
    parameters are similar: 0.125~mm vs.\ the best MLP's 0.305~mm.
\end{enumerate}

\paragraph{Known limitations carried over.}
Neither \texttt{PINN\_v2} nor \texttt{NeuralRK4} uses the per-sample
initial $z_0$; both assume $z_0 = 0$.  The training data has
$z_0 \in [0, 14000]$~mm and the magnetic field is strongly
$z$-dependent, so this is an explicit modelling approximation, flagged
in the source with \texttt{FIX NOTE} comments and documented in
\texttt{ISSUES}~\S7 and \S9.4 for resolution in gen\_3.

\subsection{Provenance and Reproducibility}
\label{sec:v2_provenance}

The fixes are traceable end-to-end:
\begin{itemize}
  \item \textbf{Code.}
    \texttt{experiments/gen\_2/models/architectures.py} --- classes
    \texttt{PINN\_v2} (added 2026-04-21, fixes A+B+C) and
    \texttt{NeuralRK4} (added 2026-04-21, fix D$'$) with the original
    \texttt{PINN} and \texttt{RK\_PINN} kept for archival comparison.
    \texttt{MODEL\_REGISTRY} now exposes \texttt{pinn\_v2} and
    \texttt{neural\_rk4}.
  \item \textbf{Training script.}
    \texttt{experiments/gen\_2/models/train.py} --- added
    \texttt{physics\_warmup\_epochs} (default 5) and a per-epoch ramp
    (Fix F).
  \item \textbf{Configs.}
    \texttt{experiments/gen\_2/configs/v2\_fixes/pinn\_v2/*.json}
    (\num{3} files, $\lambda_\text{pde} \in \{0.1, 1.0, 10.0\}$) and
    \texttt{configs/v2\_fixes/neural\_rk4/*.json} (\num{3} files,
    \texttt{tiny\_1step}, \texttt{small\_1step}, \texttt{small\_2step}).
    Each config records the list \texttt{arch\_fixes} of applied fixes
    and the \texttt{supersedes\_experiment} identifier of the legacy run
    it replaces.
  \item \textbf{Cluster.}
    \texttt{condor/submit\_v2\_fixes.sub} with
    \texttt{condor/jobs\_v2\_fixes.txt} dispatches the six v2 jobs.
    Submitted as cluster 4262459 on 2026-04-21; legacy PINN/RK\_PINN
    jobs (cluster 4262446, ProcIds 5--10) were cancelled at the same
    time.
  \item \textbf{MLflow.}
    Experiment name
    \texttt{gen\_2\_track\_extrapolation\_v2\_fixes}, stored in
    \texttt{experiments/gen\_2/mlruns/}, disjoint from the
    \texttt{gen\_2\_track\_extrapolation\_true\_gen2\_data} experiment
    that holds the MLP baselines.
\end{itemize}

% ============================================================================
\section{Companion Reports}
\label{sec:companion_reports}

This protocol is complemented by three standalone results reports stored
alongside this document in \texttt{docs/reports/}.  They are written to be
read sequentially; each references but does not duplicate the preceding
material.

\begin{itemize}
  \item \textbf{\texttt{gen1\_results.tex}} --- Initial MLP-only
    feasibility study on the V1 dataset
    (\texttt{experiments/gen\_1/V1/data\_generation/datasets/train\_50M.npz}).
    Motivates the whole project (64~kB CUDA constant-memory budget,
    Allen HLT1 integration), characterises the \texttt{twodip.rtf}
    magnetic field, presents the 17-MLP architecture sweep, and
    establishes the Pareto frontier.  Labels prefixed \texttt{gen1:}.
  \item \textbf{\texttt{gen2\_old\_data\_results.tex}} --- First PINN
    and RK\_PINN study, run with the \texttt{gen\_2} codebase on the
    V1 dataset (HTCondor cluster 4217847).  Diagnoses the four
    pathologies (P1 affine-in-$\zeta$ trunk, P2 reverse-mode AD cost,
    P3 state-averaging RK\_PINN, P4 softmax-drift of Butcher weights)
    that motivate the v2 fixes described in \S\ref{sec:v2_fixes} of
    this protocol.  Labels prefixed \texttt{gen2old:}.
  \item \textbf{\texttt{gen2\_new\_data\_results.tex}} --- v2 fixes
    (PINN\_v2, NeuralRK4) on the new variable-$\dz$ dataset
    (\texttt{train\_50M\_dz25.npz}), HTCondor clusters 4262446 (MLP
    baselines) and 4262459 (v2 fixes).  Presents the final results
    table: NeuralRK4 at 17.9k parameters reaches \SI{0.124}{\mm} position
    error and $5.0\times10^{-5}$~rad slope error, beating the best MLP
    (100.9k parameters, \SI{0.305}{\mm}) by $2.5\times$ position and
    $2\times$ slope.  Labels prefixed \texttt{gen2new:}.
\end{itemize}

% ============================================================================
\section{Summary}
\label{sec:summary}

This protocol defines a nine-configuration experiment comparing three
neural network architectures (MLP, PINN, RK\_PINN) for replacing the RK4
track extrapolator in LHCb's Allen HLT1 trigger.  The \texttt{gen\_2}
dataset extends the propagation range to $\dz \geq \SI{25}{\mm}$ to cover
VELO sensor spacing, and all MLP architectures are constrained to
$\leq$~\SI{64}{\kilo\byte} for GPU deployment.  Six confounding variables
are identified (activation, batch size, learning rate, parameter count,
training cost, and Jacobian metric sensitivity).  The analysis plan
specifies aggregate, $\dz$-binned, and momentum-binned evaluation to
address both the engineering (RQ1) and scientific (RQ2) research
questions.

% ============================================================================
% References
% ============================================================================
\begin{thebibliography}{20}

\bibitem{LHCb_detector}
R.~Aaij \textit{et al.} (LHCb Collaboration),
``The LHCb detector at the LHC,''
JINST \textbf{3}, S08005 (2008).
\href{https://doi.org/10.1088/1748-0221/3/08/S08005}{doi:10.1088/1748-0221/3/08/S08005}.

\bibitem{RK4_extrapolator}
E.~Bos and E.~"; see also
M.~Needham,
``The LHCb track extrapolation,''
LHCb-2007-140, CERN-LHCb-2007-140 (2007).
% [REF NEEDED] Exact internal note reference for TrackRungeKuttaExtrapolator.

\bibitem{LHCb_field_map}
% [REF NEEDED] Reference for twodip.rtf field map and measurement procedure.
LHCb magnetic field map, internal documentation.

\bibitem{Allen_GPU}
R.~Aaij \textit{et al.} (LHCb Collaboration),
``Allen: A high-level trigger on GPUs for LHCb,''
Comput.\ Softw.\ Big Sci.\ \textbf{4}, 7 (2020).
\href{https://doi.org/10.1007/s41781-020-00039-7}{doi:10.1007/s41781-020-00039-7}.

\bibitem{Kalman_filter}
R.~Fr\"uhwirth,
``Application of Kalman filtering to track and vertex fitting,''
Nucl.\ Instrum.\ Meth.\ A \textbf{262}, 444 (1987).
\href{https://doi.org/10.1016/0168-9002(87)90887-4}{doi:10.1016/0168-9002(87)90887-4}.

\bibitem{Glorot2010}
X.~Glorot and Y.~Bengio,
``Understanding the difficulty of training deep feedforward neural networks,''
in \textit{Proc.\ AISTATS} (2010), pp.~249--256.

\bibitem{SiLU}
S.~Elfwing, E.~Uchibe, and K.~Doya,
``Sigmoid-weighted linear units for neural network function approximation
in reinforcement learning,''
Neural Networks \textbf{107}, 3 (2018).
\href{https://doi.org/10.1016/j.neunet.2017.12.012}{doi:10.1016/j.neunet.2017.12.012}.

\bibitem{AdamW}
I.~Loshchilov and F.~Hutter,
``Decoupled weight decay regularization,''
in \textit{Proc.\ ICLR} (2019).
\href{https://arxiv.org/abs/1711.05101}{arXiv:1711.05101}.

\bibitem{Stiasny2021}
J.~Stiasny, G.~S.~Misyris, and S.~Chatzivasileiadis,
``Physics-informed neural networks for non-linear system identification
applied to power system dynamics,''
in \textit{Proc.\ IEEE PowerTech} (2021).
\href{https://doi.org/10.1109/PowerTech46648.2021.9494769}{doi:10.1109/PowerTech46648.2021.9494769}.

\bibitem{Zhai2023}
J.~Zhai, T.~Dobrowolski, and E.~Bhatt,
``Solving differential equations with physics-informed neural networks
incorporating Runge-Kutta methods,''
% [REF NEEDED] Exact journal/arXiv reference for Zhai et al. 2023.

\bibitem{Raissi2019}
M.~Raissi, P.~Perdikaris, and G.~E.~Karniadakis,
``Physics-informed neural networks: A deep learning framework for solving
forward and inverse problems involving nonlinear partial differential
equations,''
J.\ Comput.\ Phys.\ \textbf{378}, 686 (2019).
\href{https://doi.org/10.1016/j.jcp.2018.10.045}{doi:10.1016/j.jcp.2018.10.045}.

\bibitem{Wang2021}
S.~Wang, Y.~Teng, and P.~Perdikaris,
``Understanding and mitigating gradient flow pathologies in
physics-informed neural networks,''
SIAM J.\ Sci.\ Comput.\ \textbf{43}(5), A3055 (2021).
\href{https://doi.org/10.1137/20M1318043}{doi:10.1137/20M1318043}.

\bibitem{Krishnapriyan2021}
A.~S.~Krishnapriyan, A.~Gholami, S.~Zhe, R.~M.~Kirby, and
M.~W.~Mahoney,
``Characterizing possible failure modes in physics-informed neural
networks,''
in Adv.\ Neural Inf.\ Process.\ Syst.\ (NeurIPS), \textbf{34}, 26548 (2021).

\end{thebibliography}

\end{document}
